\documentclass{amsart}
\usepackage{amsmath,amssymb,amsthm,hyperref}
\newtheorem{theorem}{Theorem}
\newtheorem{lemma}{Lemma}
\newtheorem{proposition}{Proposition}
\theoremstyle{remark}
\newtheorem{remark}{Remark}
\begin{document}

\title{The uniqueness statement is equivalent to the Collatz conjecture}
\maketitle

\section*{Statement}
Let $\mathbb{D} = \{z \in \mathbb{C} : |z| < 1\}$ and $\lambda = e^{2\pi i/3}$.
We study the analytic functions $\varphi:\mathbb{D}\to\mathbb{C}$ satisfying
\begin{equation}\label{eq:FE}
3z\,\varphi(z^3) \;=\; 3z\,\varphi(z^6) \;+\; \varphi(z^2) \;+\; \lambda\,\varphi(\lambda z^2) \;+\; \lambda^2 \varphi(\lambda^2 z^2),
\qquad z\in\mathbb{D}.
\end{equation}

The question asks whether \emph{every} analytic solution has the form
$\varphi(z)=a+\dfrac{bz}{1-z}$, $a,b\in\mathbb{C}$.

\medskip
\noindent\textbf{Main result.} \emph{The proposed uniqueness statement is logically equivalent to the
Collatz ($3n+1$) conjecture.} Consequently, with current knowledge it can neither be proved nor
disproved; below we give a complete and rigorous reduction. We prove unconditionally the exact
classification of all analytic solutions, and then the equivalence with Collatz.

\section{Reduction of the functional equation to a coefficient recursion}

Since $\varphi$ is analytic on $\mathbb{D}$ it has a power series
\begin{equation}\label{eq:series}
\varphi(z)=\sum_{n\ge 0}c_n z^n,\qquad \limsup_{n\to\infty}|c_n|^{1/n}\le 1,
\end{equation}
converging absolutely on $\mathbb{D}$. All the substitutions $z\mapsto z^2,z^3,z^6,\lambda z^2,\lambda^2 z^2$
map $\mathbb{D}$ into $\mathbb{D}$, so each side of \eqref{eq:FE} is an absolutely convergent power
series on $\mathbb{D}$; by the identity theorem (uniqueness of Taylor coefficients of an analytic
function) \eqref{eq:FE} holds if and only if the coefficients of $z^k$ agree for every $k\ge 0$.

\begin{lemma}\label{lem:sym}
For every $n\ge 0$,
\[
1+\lambda^{\,n+1}+\lambda^{\,2(n+1)}=
\begin{cases}
3,& n\equiv 2 \pmod 3,\\[2pt]
0,& \text{otherwise.}
\end{cases}
\]
\end{lemma}
\begin{proof}
This is the standard root-of-unity filter: $1+\omega+\omega^2=0$ for any cube root of unity
$\omega\ne 1$, while it equals $3$ when $\omega=1$. Here $\omega=\lambda^{\,n+1}$, which equals
$1$ exactly when $3\mid n+1$, i.e.\ $n\equiv 2\pmod 3$.
\end{proof}

\begin{lemma}\label{lem:terms}
With $\varphi$ as in \eqref{eq:series},
\[
\varphi(z^2)+\lambda\,\varphi(\lambda z^2)+\lambda^2\varphi(\lambda^2 z^2)
=3\!\!\sum_{\substack{m\ge 0\\ m\equiv 2\,(3)}}\!\! c_m\,z^{2m}.
\]
\end{lemma}
\begin{proof}
The left side equals
\[
\sum_{m\ge 0} c_m\bigl(z^{2m}+\lambda(\lambda z^2)^m+\lambda^2(\lambda^2 z^2)^m\bigr)
=\sum_{m\ge0}c_m z^{2m}\bigl(1+\lambda^{m+1}+\lambda^{2(m+1)}\bigr),
\]
and we apply Lemma~\ref{lem:sym}.
\end{proof}

Substituting Lemma~\ref{lem:terms} and the expansions
$3z\,\varphi(z^3)=3\sum_{n\ge0}c_n z^{3n+1}$ and
$3z\,\varphi(z^6)=3\sum_{n\ge0}c_n z^{6n+1}$ into \eqref{eq:FE} and dividing by $3$, the
functional equation \eqref{eq:FE} is equivalent to the identity of power series
\begin{equation}\label{eq:three}
\sum_{n\ge0}c_n z^{3n+1}
=\sum_{n\ge0}c_n z^{6n+1}
+\!\!\sum_{\substack{m\ge0\\ m\equiv2\,(3)}}\!\! c_m\,z^{2m}.
\end{equation}

\begin{lemma}\label{lem:exponents}
Every monomial appearing on either side of \eqref{eq:three} has exponent $\equiv 1\pmod 3$.
Moreover, for each $n\ge 0$ the coefficient of $z^{3n+1}$ in \eqref{eq:three} yields exactly the
single relation
\begin{equation}\label{eq:recursion}
c_n \;=\;
\begin{cases}
c_{\,n/2}, & n\ \text{even},\\[2pt]
c_{\,(3n+1)/2}, & n\ \text{odd},
\end{cases}
\end{equation}
and these are \emph{all} the equations imposed by \eqref{eq:FE}.
\end{lemma}
\begin{proof}
On the left the exponents are $3n+1\equiv1\pmod3$. The first sum on the right has exponents
$6n+1\equiv1\pmod3$. In the last sum the exponents are $2m$ with $m\equiv2\pmod3$, hence
$2m\equiv 4\equiv1\pmod 3$. Thus every exponent occurring is $\equiv1\pmod3$, and matching the
coefficient of $z^{k}$ is nontrivial only for $k=3n+1$.

Fix $n\ge0$ and put $k=3n+1$. The left side of \eqref{eq:three} contributes the single term
$c_n$ (the map $n\mapsto 3n+1$ is injective). On the right:

\emph{First sum:} it contributes $c_j$ when $6j+1=3n+1$, i.e.\ $6j=3n$, i.e.\ $n=2j$. This
happens iff $n$ is even, in which case $j=n/2$ and the contribution is $c_{n/2}$ (each value
$6j+1$ occurs for a unique $j$).

\emph{Last sum:} it contributes $c_m$ when $2m=3n+1$ with $m\equiv2\pmod3$. The equation
$2m=3n+1$ has an integer solution $m$ iff $3n+1$ is even, i.e.\ iff $n$ is odd; then
$m=(3n+1)/2$. Writing $n=2t+1$ gives $m=3t+2\equiv2\pmod3$, so the congruence condition is
automatically satisfied. Hence for odd $n$ the contribution is exactly $c_{(3n+1)/2}$.

Finally, for even $n$ we have $3n+1$ odd, so the last sum does not contribute; for odd $n$ we
have $3n+1\equiv 4\pmod 6$, so the first sum (whose exponents are $\equiv1\pmod6$) does not
contribute. Therefore the coefficient identity for $z^{3n+1}$ reads precisely
\eqref{eq:recursion}. As $n$ ranges over $\mathbb{Z}_{\ge0}$, $k=3n+1$ ranges over all admissible
exponents, so \eqref{eq:recursion} for all $n\ge0$ is the complete content of \eqref{eq:FE}.
\end{proof}

\begin{remark}
For $n=0$ relation \eqref{eq:recursion} reads $c_0=c_0$, an empty constraint; so $c_0$ is a free
parameter (it will play the role of $a$).
\end{remark}

\section{The Syracuse map and the orbit structure}

Define the \emph{Syracuse map} $T:\mathbb{Z}_{\ge1}\to\mathbb{Z}_{\ge1}$ by
\[
T(n)=
\begin{cases}
n/2, & n\ \text{even},\\
(3n+1)/2, & n\ \text{odd}.
\end{cases}
\]
Note $T$ maps positive integers to positive integers, with $T(1)=2,\ T(2)=1$. By
Lemma~\ref{lem:exponents}, an analytic $\varphi=\sum c_n z^n$ solves \eqref{eq:FE} if and only if
\begin{equation}\label{eq:orbit}
c_n=c_{T(n)}\qquad\text{for every } n\ge1,
\end{equation}
with $c_0\in\mathbb{C}$ arbitrary and subject only to the analyticity growth condition
$\limsup_{n}|c_n|^{1/n}\le1$.

We analyse \eqref{eq:orbit} through the \emph{functional graph} $G$ of $T$: its vertex set is
$\mathbb{Z}_{\ge1}$ and it has a directed edge $n\to T(n)$ for every $n\ge1$. Let $\sim$ be the
equivalence relation on $\mathbb{Z}_{\ge1}$ generated by $n\sim T(n)$ (equivalently, the relation
of being in the same weakly connected component of $G$). The next lemma is the precise structural
fact that the equivalence with Collatz rests on; we make it explicit because it underlies both
directions of the main theorem.

\begin{lemma}[Weak connectivity in a functional graph]\label{lem:wc}
For $a,b\in\mathbb{Z}_{\ge1}$ we have $a\sim b$ if and only if their forward orbits eventually
meet, i.e.
\[
\exists\, i,j\ge0:\quad T^{\,i}(a)=T^{\,j}(b).
\]
\end{lemma}
\begin{proof}
Define $a\approx b$ to mean $\exists\,i,j\ge0:\ T^i(a)=T^j(b)$.

\emph{$\approx$ is an equivalence relation.} Reflexivity ($i=j=0$) and symmetry are clear. For
transitivity, suppose $T^i(a)=T^j(b)$ and $T^k(b)=T^\ell(c)$. Put $b'=T^j(b)=T^i(a)$ and
$b''=T^k(b)=T^\ell(c)$. Since $T$ is a \emph{function}, the forward orbit of $b$ is a single
sequence $b,T(b),T^2(b),\dots$; both $b'=T^j(b)$ and $b''=T^k(b)$ lie on it, so one is a forward
image of the other, say (WLOG $j\le k$) $b''=T^{\,k-j}(b')=T^{\,k-j}(T^i(a))=T^{\,(k-j)+i}(a)$.
Hence $T^{(k-j)+i}(a)=b''=T^\ell(c)$, giving $a\approx c$. (The case $k\le j$ is symmetric.)

\emph{$\sim\,\subseteq\,\approx$.} The generating relations satisfy $n\approx T(n)$ (take
$i=1,j=0$), and $\sim$ is the smallest equivalence relation containing them; since $\approx$ is an
equivalence relation containing them, $\sim\,\subseteq\,\approx$.

\emph{$\approx\,\subseteq\,\sim$.} If $T^i(a)=T^j(b)$, then $a\sim T(a)\sim\cdots\sim T^i(a)$ and
$b\sim T(b)\sim\cdots\sim T^j(b)$, and $T^i(a)=T^j(b)$; by transitivity of $\sim$ we get
$a\sim b$.
\end{proof}

\begin{lemma}\label{lem:onesim}
The Collatz conjecture (every $n\ge1$ has $T^k(n)=1$ for some $k\ge0$) holds if and only if
$\mathbb{Z}_{\ge1}$ consists of a single $\sim$-class.
\end{lemma}
\begin{proof}
($\Rightarrow$) If every $n$ reaches $1$, then for all $a,b$ we have $T^i(a)=1=T^j(b)$ for suitable
$i,j$, so $a\approx b$, hence by Lemma~\ref{lem:wc} $a\sim b$; there is one class.

($\Leftarrow$) Suppose there is a single class. Then for every $n\ge1$ we have $n\sim 1$, so by
Lemma~\ref{lem:wc} there exist $i,j$ with $T^i(n)=T^j(1)$. The forward orbit of $1$ is
$1,2,1,2,\dots$, i.e.\ $T^j(1)\in\{1,2\}$ for all $j$; and $T(2)=1$, so in either case the orbit of
$n$ reaches $1$ (if $T^i(n)=2$ then $T^{i+1}(n)=1$). Thus Collatz holds.
\end{proof}

\begin{remark}[Why divergent trajectories also give extra classes]
Lemma~\ref{lem:wc} is exactly what handles a hypothetical \emph{divergent} Collatz trajectory.
If some $n_0$ has $T^k(n_0)\to\infty$, then $n_0\not\sim 1$: otherwise Lemma~\ref{lem:wc} would
give $T^i(n_0)=T^j(1)\in\{1,2\}$ for some $i,j$, forcing the orbit of $n_0$ to reach $1$, contrary
to divergence. So a divergent trajectory produces a $\sim$-class distinct from that of $1$, just as
a nontrivial cycle does.
\end{remark}

\section{Classification of analytic solutions}

Write $\mathcal{C}$ for the set of $\sim$-classes of $\mathbb{Z}_{\ge1}$ and, for $n\ge1$, let
$C(n)\in\mathcal{C}$ be the class of $n$.

\begin{proposition}\label{prop:correspondence}
The map $\varphi\mapsto (c_n)_{n\ge0}$ (Taylor coefficients) is a bijection from the set of
analytic solutions of \eqref{eq:FE} onto the set of sequences $(c_n)_{n\ge0}$ of complex numbers
such that
\begin{itemize}
\item[(i)] $(c_n)_{n\ge1}$ is constant on each $\sim$-class, i.e.\ $c_m=c_n$ whenever
$m\sim n$ \,($c_0$ unrestricted), and
\item[(ii)] $\limsup_{n\to\infty}|c_n|^{1/n}\le1$.
\end{itemize}
In particular, \emph{every bounded} sequence $(c_n)_{n\ge0}$ satisfying (i) gives an analytic
solution, since boundedness implies (ii).
\end{proposition}
\begin{proof}
If $\varphi$ is an analytic solution, then \eqref{eq:series} gives (ii), and \eqref{eq:orbit}
states $c_n=c_{T(n)}$ for all $n\ge1$. Since $\sim$ is generated by the relations $n\sim T(n)$,
constancy along every edge $n\to T(n)$ is equivalent to constancy on every $\sim$-class; thus (i)
holds. Conversely, given a sequence satisfying (i) and (ii), the series
$\varphi(z)=\sum_{n\ge0}c_n z^n$ converges on $\mathbb{D}$ by (ii) and defines an analytic
function whose coefficients satisfy $c_n=c_{T(n)}$ (each edge $n\to T(n)$ joins two members of the
same class, on which $(c_n)$ is constant), i.e.\ \eqref{eq:orbit}, i.e.\ \eqref{eq:FE}. The two
constructions are mutually inverse because an analytic function is determined by its Taylor
coefficients.

Finally, if $(c_n)$ is bounded, say $|c_n|\le M$, then $|c_n|^{1/n}\le M^{1/n}\to1$, so (ii)
holds; hence any bounded sequence satisfying (i) yields an analytic solution.
\end{proof}

\section{Equivalence with the Collatz conjecture}

\begin{theorem}\label{thm:main}
The assertion
\begin{quote}
``every analytic $\varphi:\mathbb{D}\to\mathbb{C}$ solving \eqref{eq:FE} has the form
$\varphi(z)=a+\dfrac{bz}{1-z}$ for some $a,b\in\mathbb{C}$''
\end{quote}
is true if and only if the Collatz conjecture is true.
\end{theorem}

\begin{proof}
First we restate the target family in coefficient terms. Since
$\dfrac{bz}{1-z}=b\sum_{n\ge1}z^n$ on $\mathbb{D}$, the function $a+\dfrac{bz}{1-z}$ has Taylor
coefficients $c_0=a$ and $c_n=b$ for all $n\ge1$. Conversely, any analytic $\varphi$ with
$c_0=a$ and $c_n=b$ ($n\ge1$) equals $a+b\sum_{n\ge1}z^n=a+\dfrac{bz}{1-z}$. Hence, for an analytic
solution $\varphi$,
\begin{equation}\label{eq:star}
\varphi(z)=a+\tfrac{bz}{1-z}\ \text{for some }a,b
\quad\Longleftrightarrow\quad
c_n \text{ is the same for all } n\ge1.\tag{$\star$}
\end{equation}
(Indeed all $c_n$, $n\ge1$, being constant is exactly the condition with $b$ that common value and
$a=c_0$.)

\smallskip
($\Leftarrow$) Assume Collatz holds. By Lemma~\ref{lem:onesim} there is a single $\sim$-class, so
for any solution $\varphi$, Proposition~\ref{prop:correspondence}(i) forces $c_n=c_1$ for all
$n\ge1$. Thus $(\star)$ holds with $a=c_0$, $b=c_1$, and $\varphi(z)=a+\dfrac{bz}{1-z}$. Hence the
statement is true.

\smallskip
($\Rightarrow$) Assume Collatz fails. By Lemma~\ref{lem:onesim} there are at least two distinct
$\sim$-classes. Let $C_1=C(1)$ be the class of $1$ (it contains $\{1,2\}$) and pick a class
$C_2\ne C_1$. Define a sequence by
\[
c_0=0,\qquad c_n=\begin{cases}1,& n\in C_2,\\ 0,& n\ge1,\ n\notin C_2.\end{cases}
\]
This sequence is bounded (values in $\{0,1\}$) and constant on each $\sim$-class, so by
Proposition~\ref{prop:correspondence} the function
\[
\varphi_\star(z)=\sum_{n\in C_2} z^n
\]
is an analytic solution of \eqref{eq:FE} on $\mathbb{D}$. Now $1\in C_1\ne C_2$ gives $c_1=0$,
while any $m\in C_2$ gives $c_m=1$; since $C_2\ne\varnothing$, the coefficients $(c_n)_{n\ge1}$ are
not all equal, so by $(\star)$ the solution $\varphi_\star$ is \emph{not} of the form
$a+bz/(1-z)$. Hence the statement is false.

\smallskip
Combining the two implications proves the theorem.
\end{proof}

\begin{remark}[Honest status of the original ``prove or disprove'']
Theorem~\ref{thm:main} shows that the original question is, with present knowledge,
\emph{undecided}: a proof of the proposed uniqueness statement would prove the Collatz conjecture,
and exhibiting a single analytic counterexample (necessarily built from a nontrivial Collatz cycle
or a divergent trajectory, via the indicator $\varphi_\star$ above) would disprove Collatz. What is
established here \emph{unconditionally} is:
\begin{itemize}
\item the complete classification (Proposition~\ref{prop:correspondence}): the analytic solutions
of \eqref{eq:FE} are exactly the analytic functions whose Taylor coefficients are constant on
Syracuse $\sim$-classes, with $c_0$ free;
\item the family $a+bz/(1-z)$ are solutions, corresponding to the constant assignment
$c_n\equiv b$ ($n\ge1$), which is trivially constant on every class;
\item the exact logical equivalence (Theorem~\ref{thm:main}) of the proposed uniqueness statement
with the Collatz conjecture.
\end{itemize}
\end{remark}

\end{document}
